\section{总结与延伸}

\subsection{讲者的核心总结}

Chaofei Fan 在课程结尾总结了三点核心信息：

\begin{enumerate}
    \item \textbf{BCI 是一个令人振奋的跨学科研究方向}：它处于 AI/机器学习、神经科学和神经工程的交汇处
    \item \textbf{可用的临床系统即将到来}：语音 BCI 已从实验室演示发展到接近日常使用的水平
    \item \textbf{最重要的是}：BCI 正在为像 Howard 和 T12 这样的患者带来真实的希望——帮助他们在沉默多年后重新``发声''
\end{enumerate}

\subsection{全课知识图谱}

\begin{center}
\begin{tikzpicture}[
    node distance=0.9cm,
    every node/.style={draw, rounded corners, minimum width=3.5cm, minimum height=0.7cm, font=\small}
]
\node (neuro) {神经科学基础};
\node (record) [below=of neuro] {神经记录技术};
\node (decode) [below=of record] {神经信号解码};
\node (speech) [below=of decode] {语音 BCI 系统};
\node (future) [below=of speech] {未来方向与伦理};
\draw[->, thick] (neuro) -- (record);
\draw[->, thick] (record) -- (decode);
\draw[->, thick] (decode) -- (speech);
\draw[->, thick] (speech) -- (future);
\node[draw=none, right=1cm of neuro, text width=5.5cm, font=\scriptsize] {神经元、spike、调谐曲线、运动皮层};
\node[draw=none, right=1cm of record, text width=5.5cm, font=\scriptsize] {EEG、ECoG、微电极阵列、时空分辨率权衡};
\node[draw=none, right=1cm of decode, text width=5.5cm, font=\scriptsize] {音素解码、GRU、CTC、Beam Search、语言模型};
\node[draw=none, right=1cm of speech, text width=5.5cm, font=\scriptsize] {T12 实验、25\% WER、静默语音、实时系统};
\node[draw=none, right=1cm of future, text width=5.5cm, font=\scriptsize] {多模态 BCI、内部语音、隐私与伦理};
\end{tikzpicture}
\end{center}

\subsection{关键 Takeaways}

\begin{importantbox}{五条核心要点}
\begin{enumerate}
    \item \textbf{大脑信号可以被解码}：运动皮层的神经活动编码了丰富的运动意图信息，包括语音产生过程中的口腔面部肌肉运动
    \item \textbf{音素是解码的``中间表示''}：直接解码词汇需要海量数据，而 40 个音素的集合大大降低了数据需求，且可以组合出任意词汇
    \item \textbf{模型选择应服从问题约束}：在小数据 + 低延迟场景下，简单的 GRU + CTC 优于复杂的 Transformer encoder-decoder
    \item \textbf{语言模型是关键的``外部知识源''}：预训练语言模型弥补了神经解码器的不足，显著提高了整体系统的准确性
    \item \textbf{BCI 正在从实验走向临床}：语音 BCI 已经帮助患者在日常生活中恢复沟通能力，这项技术正在改变真实的人生
\end{enumerate}
\end{importantbox}

\subsection{拓展阅读}

\begin{itemize}
    \item Willett, F. et al. (2021). High-performance brain-to-text communication via handwriting. \textit{Nature}. \url{https://doi.org/10.1038/s41586-021-03506-2}
    \item Willett, F. et al. (2023). A high-performance speech neuroprosthesis. \textit{Nature}. \url{https://doi.org/10.1038/s41586-023-06377-x}
    \item Moses, D. A. et al. (2021). Neuroprosthesis for decoding speech in a paralyzed person with anarthria. \textit{New England Journal of Medicine}. \url{https://doi.org/10.1056/NEJMoa2027540}
    \item Metzger, S. L. et al. (2023). A high-performance neuroprosthesis for speech decoding and avatar control. \textit{Nature}. \url{https://doi.org/10.1038/s41586-023-06443-4}
    \item Card, N. S. et al. (2024). An accurate and rapidly calibrating speech neuroprosthesis. UC Davis / NPTL.
    \item Shenoy, K. V. \& Carmena, J. M. (2014). Combining decoder design and neural adaptation in brain-machine interfaces. \textit{Neuron}. \url{https://doi.org/10.1016/j.neuron.2014.08.038}
    \item BrainGate 研究联盟: \url{https://www.braingate.org/}
    \item Stanford NPTL 实验室: \url{https://nptl.stanford.edu/}
\end{itemize}

\end{document}
